\section{引言}

转座元件（transposable element，TE）是人类基因组的重要组成部分，研究者通常通过几类相互补充的表示方式研究它们。重复序列家族资源描述分类体系和共识序列模型，基因组注释记录单个或代表性位点，表达数据集测量特定样本中的转录水平，而生物医学文献则报道 TE 与基因、疾病及其他实体之间的关联。这些表示方式不能相互替代。家族共识序列并不等同于每个基因组拷贝的序列，观察到某个基因组位点并不能证明其具有活性，文献报道的关联本身也不能建立因果关系。因此，对 TE 证据的解释同时取决于证据来源和观察单位。

现有专业资源在特定证据类型上具有较大深度。Repbase 和 Dfam 支持重复序列家族分类与参考序列模型 \citep{Bao2015Repbase,Kojima2018HumanRepbase,Wheeler2013Dfam}。HERVd、dbRIP、euL1db 和 TranspoGene 分别整理内源性逆转录病毒序列、逆转录转座子插入多态性、样本特异性 L1HS 插入或基因环境注释 \citep{Paces2002HERVd,Wang2006dbRIP,Mir2015euL1db,Levy2008TranspoGene}。较新的资源进一步关注疾病或表达证据。HervD Atlas 整理 HERV 与疾病的关联并通过交互式知识图谱提供访问，CancerHERVdb 聚焦癌症中的 HERV 表达，TE-SCALE 则提供人类癌症单细胞水平的 TE 表达和 TE-基因共表达信息 \citep{Li2024HervDAtlas,Stricker2023CancerHERVdb,Meng2026TESCALE}。这些专业资源能够以较高分辨率回答重要问题，通用整合资源不应声称能够取代它们。

尚未充分解决的实际问题是跨证据层导航。研究者从一个 TE 名称出发，可能需要确定其分类、查看参考序列或基因组记录、比较不同表达情境、识别相关基因，并打开支持某条疾病或基因关系的论文。当这些步骤分散在不同资源中时，名称和证据边界需要反复对齐。相反，如果把异质记录合并到一个不加区分的图中，较弱或间接的证据又可能显得比实际更强。因此，有价值的整合不仅要连接记录，还必须保留其来源和解释限制。

本文介绍 TE-KG，一个以人类 TE 为中心的资源，用于连接文献提取的生物医学关系、分类体系、代表性基因组与序列记录、表达谱以及特定情境下的共表达网络。TE-KG 通过图、路径、表格、可视化、下载和自然语言等入口访问这些证据层。其核心设计原则是：整合应当使不同证据能够被联合访问，但不能抹去各来源能够支持的结论边界。本文报告当前数据库内容、数据处理与运行时架构以及主要科学访问流程，并明确限制生物学解释的因素和公开发布前仍需补充的信息。
